\documentclass[11pt]{article}

\usepackage[a4paper,margin=1in]{geometry}
\usepackage{setspace}
\usepackage{parskip}
\usepackage{titlesec}

\setstretch{1.1}

\title{Daughters, Labels, and the Promise of Long-Term Collaboration}
\author{}
\date{}

\begin{document}

\maketitle

Today I find myself thinking about my daughters --- about the world they are growing up in, and about the world we are, collectively, building for them. It is difficult not to feel the weight of history and the volatility of the present. Yet it is precisely in that tension that hope becomes meaningful.

For me, the Theory of Long-Term Collaboration (TLC) represents such hope. It offers a different lens through which to understand how we relate to one another and how societies might organize themselves more wisely.

At the heart of TLC lies a simple but transformative idea: collaboration is not merely interaction between individuals, groups, or identities. It is interaction between learning networks. A learning network is defined not by its outer shell, but by its capacity to learn --- by how it processes information, integrates feedback, updates its models of the world, and adapts its behavior.

From this perspective, the surface characteristics that so often dominate social categorization --- gender, skin color, physical ability, language, inherited religious or cultural labels --- are secondary. They describe context, not essence. A learning network housed in a female body is still a learning network. A learning network that speaks Hebrew, Dutch, Arabic, or English is still a learning network. What truly matters is the quality of its learning: what information it receives, how accurate that information is, whether feedback is allowed to circulate, and whether communication remains open.

This shift has profound implications. If collaboration is the interaction between learning networks, then the integrity of what passes between them becomes critical. False information, withheld evidence, and denied realities do not merely constitute moral shortcomings. They actively degrade the learning capacity of the system. They distort feedback loops. They impair adaptation. In short, they are structurally destructive.

Seeing society as a web of learning networks also reframes our understanding of norms and rules. Lying and denial cease to be wrong only because they violate ethical codes; they become harmful because they undermine collective learning itself. At the same time, injustices such as racism or patriarchy can be understood not only as moral failures but as systemic distortions that restrict whose data counts, whose feedback is heard, and whose learning is allowed to shape the shared model of reality.

This foundation feels more universal. It does not rely on a specific tradition, ideology, or identity group. It rests on something more fundamental: if we are interdependent learning systems, then our survival and flourishing depend on protecting the conditions that allow accurate learning across the network.

And this brings me back to my daughters. They are still culturally labeled as Jewish. History reminds us that in another time and place, that label alone could have determined their fate. The tragedy of such moments lies precisely in reducing complex learning beings to inherited categories.

TLC offers an antidote to that reduction. Within its frame, my daughters are not defined by a label attached at birth. They are defined by their capacities to learn, to update, to collaborate, to contribute. And crucially, they would be expected to extend that same recognition to others. The symmetry matters. Neutrality is not indifference; it is the disciplined refusal to collapse a person into a stereotype.

If enough of us were to adopt this orientation, old labels would gradually lose their grip. They would not disappear from history or culture, but they would cease to function as primary determinants of worth or belonging. What would matter instead is whether we contribute to the shared learning of the whole.

That thought makes me genuinely hopeful. Not only for my daughters, but for all of us.

\end{document}